\chapter{Происхождение семейного действия на multiplicity-среде}
\label{ch:v8-baryon-c0-multiplicity-environment-so3-action-parent-origin-gate}

Три текущих коннектора удобно записать как матричные единицы в
$\operatorname{Hom}(\mathbb R^2,\mathbb R^3)$:
\begin{equation}
 V_0=E_{00},\qquad V_1=E_{11},\qquad V_2=E_{21},\qquad
 W_{\rm cur}=\operatorname{span}_{\mathbb R}\{V_0,V_1,V_2\},
 \quad \dim W_{\rm cur}=3.
 \label{eq:v8-baryon-c0-current-arrow-matrix-model}
\end{equation}
Здесь строки нумеруют три старых endpoint-направления, а столбцы ---
два source-состояния $s_0,a_0$. Поэтому multiplicity-индекс не является
отдельным тензорным множителем: первая стрелка использует другой source,
чем вторая и третья.

На target действует стандартная тройка с генераторами
\begin{equation}
 J_1=\begin{pmatrix}0&0&0\\0&0&-1\\0&1&0\end{pmatrix},\quad
 J_2=\begin{pmatrix}0&0&1\\0&0&0\\-1&0&0\end{pmatrix},\quad
 J_3=\begin{pmatrix}0&-1&0\\1&0&0\\0&0&0\end{pmatrix},
 \qquad [J_1,J_2]=J_3.
 \label{eq:v8-baryon-c0-arrow-target-so3-generators}
\end{equation}
Не существует компенсирующего нетривиального ортогонального действия на
двумерном source: всякий гомоморфизм $\mathfrak{so}(3)\to\mathfrak{so}(2)$
нулевой, поскольку target абелев, а $\mathfrak{so}(3)$ совершенна:
\begin{equation}
 \rho_{\rm src}:SO(3)\longrightarrow O(2)
 \quad\Longrightarrow\quad d\rho_{\rm src}=0,
 \qquad \rank\mathcal C_{\rm src}=3.
 \label{eq:v8-baryon-c0-source-two-dimensional-so3-triviality}
\end{equation}

Следовательно, инфинитезимальное действие на стрелках имеет вид
\begin{equation}
 \delta_a(V)=J_aV.
 \label{eq:v8-baryon-c0-arrow-infinitesimal-so3-action}
\end{equation}
Прямой точный тест показывает, что четыре из девяти образов выходят из
текущего подпространства:
\begin{equation}
 \#\{(a,i):J_aV_i\notin W_{\rm cur}\}=4,
 \qquad J_3E_{00}=E_{10}\notin W_{\rm cur}.
 \label{eq:v8-baryon-c0-current-arrow-space-not-so3-invariant}
\end{equation}
Таким образом, старое семейное действие нельзя просто ограничить на три
имеющиеся стрелки.

\section{Минимальное ковариантное замыкание}

Одношаговое $SO(3)$-замыкание уже порождает всё пространство стрелок:
\begin{equation}
 \operatorname{span}\bigl(W_{\rm cur},J_aW_{\rm cur}\bigr)
 =\operatorname{Hom}(\mathbb R^2,\mathbb R^3),
 \qquad \rank=6.
 \label{eq:v8-baryon-c0-minimal-so3-arrow-closure-rank}
\end{equation}
Точно недостающие матричные единицы равны
\begin{equation}
 W_{\rm miss}=\operatorname{span}\{E_{10},E_{20},E_{01}\},
 \qquad W_{\rm cur}+W_{\rm miss}=M_{3\times2}(\mathbb R).
 \label{eq:v8-baryon-c0-missing-so3-connector-arrows}
\end{equation}
Это три новые комплексные стрелки, то есть шесть Hermitian-квадратур.
Их точный матричный ранг равен двенадцати вместе с исходными шестью
квадратурами, а условная размерность полного кадра становится
\begin{equation}
 \rank\mathcal F_{\rm connector}^{SO(3)}=12,
 \qquad \dim\mathcal F_{\rm full}^{\rm cond}=51+6=57.
 \label{eq:v8-baryon-c0-so3-completed-frame-rank}
\end{equation}
Алгебра $M_5$ содержит эти матричные единицы после замыкания произведениями,
но текущий шумовой кадр их не содержит как независимые скачки. Поэтому
алгебраическое наличие не является происхождением ковариантного frame.

При тривиальном source-действии замкнутое пространство раскладывается как
\begin{equation}
 \operatorname{Hom}(\mathbb R^2,\mathbb R^3)
 \simeq\mathbf3\otimes\mathbf1^{\oplus2}
 \simeq\mathbf3\oplus\mathbf3.
 \label{eq:v8-baryon-c0-so3-closed-arrow-representation-decomposition}
\end{equation}
Поэтому семейная тройка входит уже с кратностью два:
\begin{equation}
 \dim\operatorname{Hom}_{SO(3)}
 \bigl(\mathbf3,\mathbf3\oplus\mathbf3\bigr)=2,
 \qquad \rank\mathcal C_{\rm Hom}=16\ \text{на}\ 18\ \text{неизвестных}.
 \label{eq:v8-baryon-c0-so3-closed-family-hom-dimension}
\end{equation}
Общий интертвейнер имеет вид
\begin{equation}
 T_{u,v}=\begin{pmatrix}uI_3\\vI_3\end{pmatrix},
 \qquad T_{u,v}^{\mathsf T}T_{u,v}=(u^2+v^2)I_3.
 \label{eq:v8-baryon-c0-so3-closed-general-intertwiner}
\end{equation}
После изометрической нормировки остаётся не единственная карта, а
проективная source-линия:
\begin{equation}
 u^2+v^2=1,\qquad (u,v)\sim(-u,-v),\qquad
 [u:v]\in\mathbb {RP}^1.
 \label{eq:v8-baryon-c0-so3-closed-source-line-orbit}
\end{equation}
Значит, минимальное ковариантное расширение уменьшает неоднозначность
$\mathbb {RP}^2$ до $\mathbb {RP}^1$, но не выбирает карту $c_0$.

Итоговый parent-origin реестр равен
\begin{equation}
 N_{SO(3)\text{-origin}}=0/5:
 \quad\{H_{21}\text{-restriction},\ endpoint,\ W_{\rm cur},\
 M_5\text{-frame},\ source\ line\}.
 \label{eq:v8-baryon-c0-so3-action-parent-origin-ledger}
\end{equation}

\begin{theorem}[Удвоение при минимальном семейном замыкании]
Текущее трёхмерное пространство коннекторов не инвариантно относительно
стандартного $SO(3)_{\rm fam}$. Его минимальное ковариантное замыкание равно
всему шестимерному $M_{3\times2}(\mathbb R)$ и содержит две стандартные
тройки. Поэтому даже после добавления трёх недостающих комплексных стрелок
семейный интертвейнер остаётся семейством, параметризованным
$\mathbb {RP}^1$.
\end{theorem}

\begin{nogocriterion}[Алгебраическое замыкание не выбирает frame]
Наличие недостающих матричных единиц в полной алгебре $M_5$ не означает,
что они входят в GKSL-кадр с каноническими коэффициентами. Без явного
parent-источника шести стрелок и выбора source-линии нельзя объявлять
$SO(3)$-действие или одиночную карту $c_0$ выведенными.
\end{nogocriterion}

\begin{protocol}\begin{enumerate}
\item текущая стрелочная модель: \textbf{$\operatorname{span}\{E_{00},E_{11},E_{21}\}$};
\item source-представление $SO(3)$ в размерности два: \textbf{тривиально};
\item текущая тройка инвариантна: \textbf{нет};
\item число инфинитезимальных выходов: \textbf{$4$};
\item минимальная размерность замыкания: \textbf{$6$};
\item недостающих комплексных стрелок: \textbf{$3$};
\item условный полный frame: \textbf{$57$};
\item тип замкнутой среды: \textbf{$\mathbf3\oplus\mathbf3$};
\item размерность family-Hom: \textbf{$2$};
\item остаточный селектор: \textbf{$\mathbb {RP}^1$};
\item выведенных parent-источников: \textbf{$0/5$};
\item следующий гейт: \textbf{селектор source-линии $SO(3)$-замкнутой среды}.
\end{enumerate}\end{protocol}

Машинный сертификат сохранён в
\path{s2t_v8_baryon_c0_multiplicity_environment_so3_action_parent_origin_gate_results.json}.